\documentclass[11pt]{article}
\usepackage{amsmath,amssymb,amsthm}
\usepackage{booktabs}
\usepackage{hyperref}
\usepackage{geometry}

\usepackage[ruled,vlined,norelsize]{algorithm2e}
\usepackage{xcolor}
\newcommand{\abs}[1]{\lvert#1\rvert}
\newcommand{\Fq}{\mathbb{F}_q}
\newcommand{\Pperp}{P^{\perp}}
\newcommand{\rk}{\operatorname{rank}}
\newcommand{\GRS}{\mathrm{GRS}}

\geometry{margin=1in}

\newtheorem{theorem}{Theorem}
\newtheorem{lemma}{Lemma}
\newtheorem{proposition}{Proposition}
\newtheorem{corollary}{Corollary}
\newtheorem{conjecture}{Conjecture}
\newtheorem{heuristic}{Heuristic}
\theoremstyle{definition}
\newtheorem{definition}{Definition}
\theoremstyle{remark}
\newtheorem{remark}{Remark}

\title{Key Recovery Against RLCE via the GIJS Distinguisher}
\author{Jacob Alperin-Sheriff\thanks{The author was a member of the NIST PQC
    standardization team during Round 1 and participated in the evaluation of
    RLCE.}\\{\small Independent Researcher}\\{\small
\texttt{jacobmas@gmail.com}}}
\date{\today}

\begin{document}
\maketitle

\begin{abstract}
    We present an experimental cryptanalysis of the RLCE
    cryptosystem~\cite{DBLP:journals/corr/Wang15aa, RLCE-spec}
    at the heretofore secure $w=n-k$ parameter regime.  Our attack relies
    crucially on the  recent Ghoshal--Ishai--Jain--Sun (GIJS)
    distinguisher~\cite{GIJS26}. The attack shortens the public key at
    random positions, runs the GIJS distinguisher with \emph{every}
    column as hold-out, and uses this distinguisher to identify twin pairs based on how they appear across trials.
    Three pairs so identified are unmixed by a root-finding cascade,
    and Sidelnikov--Shestakov on the resulting $k \times (k+3)$ pure
    matrix recovers the key.

    We experimentally verify our attack by breaking the ID-0 (128-bit) in 30.5
    core-hours (about 3 hours of wall time) on a standard higher-end laptop.
\end{abstract}


\section{Introduction}

The RLCE cryptosystem of Wang~\cite{DBLP:journals/corr/Wang15aa, RLCE-spec} was one of the
69 schemes submitted to NIST's postquantum cryptography standardization
process. RLCE builds its public key from a generalized Reed--Solomon code
$GRS_{n,k}$ by injecting $w$ \emph{twin pairs}: for each pair, two GRS columns
are mixed with two random columns via a $2 \times 2$ invertible matrix, and all
$N = n + w$ columns are permuted. The three unbroken parameter crucially all
have $w = n-k$, the exact boundary where the CLT square distinguisher provably
fails.

This scheme was one of many~\cite{DBLP:conf/pqcrypt/Wieschebrink09,
      DBLP:journals/dcc/BaldiBC13, DBLP:conf/wcc/Wieschebrink16,
      DBLP:journals/iet-it/AlmeidaBNS24, DBLP:conf/redundancy/Krouk20} attempts to
patch Niederreiter's original attempt to instantiate the
McEliece~\cite{McEliece1978} cryptosystem with GRS
codes~\cite{DBLP:journals/dc/Niederreiter86}, which was itself broken by the
Sidenikov--Shestakov attack~\cite{sidelnikov1992insecurity}. The original
scheme and all of these other previous attempts have fallen to cryptanalysis.

Due to security concerns raised by the Couvreur--Lequesne--Tillich
(CLT)~\cite{CLT18} attack on smaller parameters, NIST did not advance RLCE to
the second round; it appears among the schemes for which NIST ``lacked full
confidence in security''~\cite{Moody19}. See~\cite{NIST8240} for the Round~1
report. No subsequent public attack has been advanced against the $w = n-k$
parameter sets since, and due to the large parameter sizes the scheme may not
have received any attention at all in the intervening period. However, with the
recent improvements in AI cryptanalysis ~\cite{cryptoeprint:2026/1593,
      cryptoeprint:2026/2266} causing concerns from eminent cryptographers that key
exchange/key encapsulation itself may be at risk ~\cite{green2026pkc}, it is
likely important to subject even large parameter schemes to renewed
cryptanalysis attempts in order to enable a potential broadening of
standardized schemes for cryptographic agility as a hedge against AI
breakthroughs against lattices, isogenies or quasi-cyclic codes.

There has been recent hubbub over the Ghoshal--Ishai--Jain--Sun~\cite{GIJS26}
provable distinguisher and heuristic key recovery for the venerable and ancient
(by cryptographic standards) McEliece scheme~\cite{McEliece1978} and subsequent
improvement papers~\cite{cryptoeprint:2026/1984}. Most notably it has led to
the German Federal Office for Information Security (BSI) shifting its guidance
from wholesale approval to recommending against the use of the Classic McEliece
scheme in new deployments~\cite{bsi2026classicmceliece}. In the course of a
very long thread on NIST's PQC forum over exactly how devastating or not these
attacks are, in practice, Dr. Wang commented that "I do not think these attacks
on McEliece will work on RLCE. I still believe RLCE has sufficient randomness
in the generator matrix and is a good alternative for McEliece
schemes,"~\cite{wang2026pqcforum} which piqued the interest of this work's
author.

We contradict Wang's claim for \emph{both} distinguishing and key recovery by
using the GIJS distinguisher: the distinguishing step is proved, and the
key-recovery step is an experimentally verified cryptanalysis~\cite{GIJS26}.

\subsection{The Attack}\label{sec:intro-attack}

In slightly more detail, our attack proceeds as follows.
\begin{enumerate}
      \item \textbf{Shorten} the public key at $\ell = k - k'$ random positions,
            for
            some \emph{small} $k'$.
      \item \textbf{Distinguish} using the recent
            Ghoshal--Ishai--Jain--Sun~\cite{GIJS26}
            distinguisher (Theorem~\ref{thm:ghoshal}), which operates beyond the
            Schur-square barrier, applied \emph{directly} with every column as
            hold-out.
      \item \textbf{Group} columns across trials via co-occurrence with twin pairs.
      \item \textbf{Recover} the original GRS columns from 3 twin pairs.
      \item \textbf{Solve} using Sidelnikov-Shestakov.
\end{enumerate}

A key insight is that GIJS will still work even with some random columns, and
that, furthermore, the pairs of noisy columns provide useful information to the
GIJS distinguisher. We show in Section~\ref{sec:gijs} that the GIJS
distinguisher for GRS subcodes should work on shortened RLCE public keys
whenever $sC + U > dD$, where $s$ is the jet parameter, $d$ is the degree of the
multivariate polynomials, $C$ is the number of clean (either from shortening at
one of a twin pair's columns or from an initially clean column) columns in the
shortened code, and $U$ is the number of twin pairs remaining in the shortened
code.

Extensive experimental testing shows that the GRS distinguisher on the
parameter sets of interest essentially never fails to diagnose a ``clean" GRS
column as GRS and a noisy column as ``random." This leads to practical attacks.

In short, we simply choose uniformly at random columns to shorten at for a
relatively small (less than 100 for the sets considered) number of separate
\emph{trials}. In each trial, we run the GIJS distinguisher on the shortened
code, with each column held out separately as the holdout column. This enables
us to learn which columns in the shortening remain part of a twin pair (meaning
both it and its pair column are in the shortened code) in that shortening, and
refer to these columns as \emph{flagged}.

To recover twin pairs, we observe that (assuming no errors in the
distinguisher) a column is \emph{flagged} in a trial if and only its twin pair
column is also \emph{flagged} in that trial. This fact means that after a
relatively small number of randomized trials, for a given column A, there will
be only other column B that is flagged in and only in trials in which column A
is flagged, and that this column B is the twin pair for column A.

After recovering all twin pairs, we use some clever ideas in\ref{sec:fastp2} to
unmix 3 of these twin pairs back into actual GRS columns. We note that a
simpler route to do so would involve simply trying all $(q-1)^3$ non-degenerate
linear combinations of the pairs, one of which will correctly unmix all 3.

Finally, after we have recovered a $k \times (k+3)$ pure GRS matrix, we are able to
run the celebrated~\cite{sidelnikov1992insecurity} to recover the code
structure (the scrambling matrix $S$). At this point, unmixing the remainder of the twin pairs
becomes trivial (see~\cite{CLT18}). We end up with the full private key.

\subsection{Experimental Success}\label{sec:exp-success}

We were able to run a full key recovery (see Section~\ref{sec:total-costs}) at
the 128 bit parameters (ID-0) in the specification that took approximately 30.5
hours of core work on a standard high-end laptop to recover the key. We have
also given key recovery data for several other smaller parameter sets for
comparison in the $w=n-k$ regime; see Section~\ref{sec:total-costs} for
details.

While we did not execute an experimental attack on $ID-1$ and $ID-2$ due to a
lack of desire to spend personal funds, we estimate that a cloud based attack should be possible
for no more than a few hundred dollars; the biggest issues, as mentioned
in~\cite{GIJS26}, are the memory costs of the attack. However, these are easily
solvable in the cloud for the given parameter sets.

\subsection{Further Work}\label{sec:further-work}
The scheme is completely broken in the $w=n-k$ parameter regime, as a back of the envelope 
computation estimates that a public key would need to be in the petabytes with hundreds of megabytes per 
ciphertext in order to reach $128$ bits of security. However, there is a natural
defensive option; push the number of twin pairs $w$ to its maximum
$w=n$, eliminating pure columns entirely ($N=2n$). Our attack does not work in
this regime when $2k \geq n$, so one could potentially propose this as a patch.
However, due to the matrix becoming wider while the error correction regime
does not change, ISD attack perform better, necessitating a bump in $k$, and
significantly raising parameter sizes.

Below, we have calculated parameter sets that seem to be immune to the attacks
given in this paper, as well as to ISD-based attacks. 

\begin{table}[!ht]
      \caption{Proposed $w=n$ parameter sets immune to our attack and ISD}
      \begin{tabular}{lrrrrrrrrr}
            \toprule
                      & $n$  & $k$  & $t$ & $N$  & field    & Stern & pk (KB)   & ct (kB)\\
            \midrule
            128-fixed & 1191 & 619  & 286 & 2382 & $2^{11}$ & 140.2 & 1980    & 3.2 \\
            192-fixed & 1774 & 922  & 426 & 3548 & $2^{11}$ & 204.0 & 4393    & 4.8 \\
            256-fixed & 2355 & 1225 & 565 & 4710 & $2^{12}$ & 268.1 & 8452    & 6.9 \\
            \bottomrule
      \end{tabular}
\end{table}

We do not claim that these parameters are secure, and if forced, we would bet
against their security given the long, long history of attacks on GRS-McEliece
schemes and various attempts to patch. We merely give them as new targets and
the smallest scheme parameters secure given our new attacks for future research. 

\section{Preliminaries}

\begin{definition}[Generalized Reed--Solomon code]
    Let $q$ be a prime power, $n \leq q$, $1 \leq k \leq n$. Choose distinct
    $x_1,\dots,x_n \in \mathbb{F}_q$ (evaluation points) and nonzero
    $y_1,\dots,y_n
        \in \mathbb{F}_q^*$ (column multipliers). The \emph{generalized
        Reed--Solomon
        code} $\mathrm{GRS}_{n,k}(x,y)$ is the $[n,k]$ linear code with
    generator
    matrix
    \[
        (G_{\mathrm{GRS}})_{i,j} = y_j x_j^i, \qquad i = 0,\dots,k-1,\; j =
        1,\dots,n,
    \]
    i.e., the evaluation of all polynomials of degree $< k$ at the $x_j$,
    scaled by
    $y_j$. GRS codes are MDS: minimum distance $d = n-k+1$.
    See~\cite[Ch.~10]{MS77}
    or~\cite[Ch.~5]{HP03}.
\end{definition}

\begin{definition}[Shortening]
    Let $C$ be an $[n,k]$ code and $L \subset \{1,\dots,n\}$. The
    \emph{shortening}
    of $C$ at $L$, denoted $C_L$ or $\mathcal{S}_L(C)$, is obtained by
    restricting to
    codewords that vanish on $L$, then deleting those positions. If the
    generator
    columns at $L$ are linearly independent, $\dim C_L = k - |L|$ and the
    length is
    $n - |L|$.
\end{definition}

\subsection{The RLCE scheme}

We use the $S \cdot G_s \cdot P_1 \cdot A \cdot P_2$ formulation from the RLCE
specification~\cite{Wan16,RLCE-spec}. Everything happens over the finite field $\Fq$ for 
$q=2^m$. 

Let $G_s$ be a $k \times n$ GRS generator
matrix. Key generation:
\begin{itemize}
    \item $S$: uniform random $k \times k$ invertible matrix (scrambler).
    \item $P_1$: uniform random $n \times n$ permutation matrix.
    \item $A$: $n \times (n+w)$ \emph{expansion} matrix. For $w$ designated
          positions $j_1,\dots,j_w$, column $j_t$ of $G_s P_1$ is replaced by
          two columns
          \[
              a_t g_{j_t} + c_t r_t, \qquad b_t g_{j_t} + d_t r_t,
          \]
          where $r_t \in \mathbb{F}_q^k$ is uniform random and
          $\bigl(\begin{smallmatrix}a_t&b_t\\c_t&d_t\end{smallmatrix}\bigr) \in
              \mathrm{GL}_2(\mathbb{F}_q)$ is uniform random. All other columns
          pass through
          unchanged. The pairs $(a_t g_{j_t} + c_t r_t,\; b_t g_{j_t} + d_t
              r_t)$ are the
          \emph{twin pairs}.
    \item $P_2$: uniform random $(n+w) \times (n+w)$ permutation matrix.
\end{itemize}
The public key is $G_{\mathrm{pub}} = S \cdot G_s \cdot P_1 \cdot A \cdot P_2$,
a $k \times (n+w)$ matrix. The secret key is $(S, G_s, P_1, A, P_2)$.
Encryption is $c = m G_{\mathrm{pub}} + e$ with $\mathrm{wt}(e) \leq t$;
decryption unpermutes by $P_2^{-1}$, unmixes each twin pair via the $2 \times
    2$ inverse, unpermutes by $P_1^{-1}$, and decodes the GRS code. Security of
the post-CLT paper updated version
relies
on the indistinguishability of $G_{\mathrm{pub}}$ from random solely at $w =
    n-k$; for a more precise statement of
security, consult the specification.


For completeness, we reproduce the specification's parameter sets below.

\begin{table}[htbp]
\centering
\footnotesize
\begin{tabular}{lccccc}
\toprule
Set &  n & k & t & w & m \\
\midrule
ID-0 & 630 & 470 & 80 & 160 & 10 \\
ID-1 & 1000 & 764 & 118 & 236 & 10 \\
ID-2 & 1360 & 800 & 280 & 560 & 11 \\
\bottomrule
\end{tabular}
\caption{Parameter sets from the RLCE specification}
\label{tab:spec-params}
\end{table}
\subsection{The CLT Attack}

We work with $G_{\mathrm{GRS}}$ in canonical form as above. The RLCE public key
has $N = n + w$ columns: $n - w$ pure GRS columns and $w$ twin pairs, each
$(c_i, c_{\tau(i)})$ of the form
\[
    c_i = a G_j + c R, \qquad c_{\tau(i)} = b G_j + d R,
\]
with $G_j$ a GRS column, $R \in \mathbb{F}_q^k$ random, and
$\bigl(\begin{smallmatrix}a&b\\c&d\end{smallmatrix}\bigr)$ invertible, all
permuted by secret $\pi$.

Shortening at $L$ deletes those positions after restricting to codewords
vanishing there. By~\cite[Lemma~24]{CLT18}, shortening at one twin-pair member
\emph{derandomizes} the other. If $L$ hits every pair, $C_L = M \cdot
    G'_{\mathrm{GRS}}$ with $M$ generic $k' \times k_{\mathrm{GRS}}$, $k' = k -
    |L|$. If $L$ misses $j$ pairs, those $2j$ columns remain mixed (noisy).

The ambient GRS dimension is $k_{\mathrm{GRS}} = k' + w - j$ for $j$ unhit
pairs.

\paragraph{High-level overview.} The Couvreur--Lequesne--Tillich
attack~\cite{CLT18} proceeds in three phases. First, \emph{distinguish}: the
Schur square $C^{\star 2}$ has abnormally low dimension when the gap $c = (n-k)
    - w$ exceeds $(3+\sqrt{16w+1})/2$, revealing twin-pair positions. Second,
\emph{unmix}: for each identified pair, the $2 \times 2$ mixing is inverted to
recover pure GRS columns. Third, \emph{recover}:
Sidelnikov--Shestakov~\cite{sidelnikov1992insecurity} on the
reconstructed GRS code yields the secret support and multipliers. The attack
\emph{succeeds} for $w$ sufficiently below $n-k$ but \emph{fails} at the
shipped $w = n-k$ boundary, where the square is provably full and no
distinguisher was known.

\subsection{Hasse Derivatives, Multiplicity Codes, and Jets}\label{sec:hasse}

We follow the treatment of~\cite{GIJS26}. For $\mathbf{a} = (a_1,\dots,a_{k'})
    \in \mathbb{Z}_{\geq 0}^{k'}$ write $|\mathbf{a}| = \sum_i a_i$ for its
$\ell_1$-norm, and let
\[
    A_{<s,k'} := \{\mathbf{a} \in \mathbb{Z}_{\geq 0}^{k'} : |\mathbf{a}| < s\}
\]
be the set of multi-indices of order below $s$.  We write
$\mathbf{X}^{\mathbf{a}} := \prod_i X_i^{a_i}$ for the corresponding monomial.

\begin{definition}[Hasse derivatives~{\cite[Def.~2.7]{GIJS26}}]
    Let $F \in \mathbb{F}[X_1,\dots,X_{k'}]$ over a field $\mathbb{F}$.  The
    \emph{Hasse derivative} of $F$ with respect to
    $\mathbf{a} \in \mathbb{Z}_{\geq 0}^{k'}$ is
    \[
        \partial^{[\mathbf{a}]} F
        = \sum_{\mathbf{b} \geq \mathbf{a}}
        \binom{\mathbf{b}}{\mathbf{a}} \cdot \mathrm{Coeff}_{\mathbf{b}}(F)\,
        \mathbf{X}^{\mathbf{b}-\mathbf{a}},
    \]
    where $\mathrm{Coeff}_{\mathbf{b}}(F)$ is the coefficient of
    $\mathbf{X}^{\mathbf{b}}$ in $F$ and
    $\binom{\mathbf{b}}{\mathbf{a}} = \prod_i \binom{b_i}{a_i}$.  Equivalently,
    Hasse derivatives are the finite-field analogue of ordinary partial
    derivatives, characterized by the formal Taylor expansion
    \[
        F(\mathbf{X}+\mathbf{Z})
        = \sum_{\mathbf{a}} (\partial^{[\mathbf{a}]}
        F)(\mathbf{X})\,\mathbf{Z}^{\mathbf{a}}.
    \]
    For a univariate polynomial $f(Z) = \sum_{j=0}^{d} c_j Z^j$ we write
    $\partial^{[r]} f(Z) = \sum_{j=r}^{d} \binom{j}{r} c_j Z^{j-r}$; in
    characteristic $0$ this equals $f^{(r)}(Z)/r!$.  The Hasse-derivative chain
    rule holds in any characteristic, which is what makes multiplicities usable
    over finite fields.
\end{definition}

\begin{definition}[Hasse jet~{\cite[Def.~2.9]{GIJS26}}]
    For $\mathbf{x} \in \mathbb{F}_q^{k'}$ and $s \geq 1$, the \emph{Hasse jet
        of
        order $<s$} of $F$ at $\mathbf{x}$ (also the \emph{order-$s$
        evaluation} of
    $F$ at $\mathbf{x}$) is the vector of all Hasse derivatives of order below
    $s$ evaluated at $\mathbf{x}$:
    \[
        J_{\mathbf{x}}^{<s}(F) := \bigl((\partial^{[\mathbf{a}]}
        F)(\mathbf{x})\bigr)_{\mathbf{a} \in A_{<s,k'}}.
    \]
    We say $F$ \emph{vanishes with multiplicity} (at least) $s$ at $\mathbf{x}$
    if $J_{\mathbf{x}}^{<s}(F) = 0$.
\end{definition}

\begin{definition}[Multiplicity code~{\cite[Def.~2.10]{GIJS26},
    \cite{DBLP:journals/jacm/KoppartySY14}}]
    The multiplicity code of order-$s$ evaluations of polynomials of total
    degree
    at most $d$ in $k'$ variables over $\mathbb{F}_q$ has length $q^{k'}$ over
    the
    alphabet $\Sigma = \mathbb{F}_q^{\binom{k'+s-1}{s-1}}$, indexed by
    $\mathbf{x} \in \mathbb{F}_q^{k'}$, with codeword
    $\mathrm{Encode}_{s,d,k',q}(F) = (J_{\mathbf{x}}^{<s}(F))_{\mathbf{x} \in
        \mathbb{F}_q^{k'}}$ for each $F$ of total degree at most $d$.
\end{definition}

The key analytic fact is Hermite interpolation: over any field, a univariate
polynomial of degree at most $d$ is determined by its order-$<s$ Hasse jets at
enough points.

\begin{lemma}[Hermite interpolation of Hasse
    derivatives~{\cite[Lem.~2.8]{GIJS26}}]
    Let $f$ be a univariate polynomial of degree at most $d$ over a finite
    field
    $\mathbb{F}$, let $\alpha_1,\dots,\alpha_h$ be distinct elements of
    $\mathbb{F}$, and let $\mu_1,\dots,\mu_h$ be positive integers with
    $\sum_i \mu_i > d$.  Then $f$ is recovered in
    $\mathrm{poly}(d,h,\log |\mathbb{F}|)$ time from the values
    $(\partial^{[j]} f)(\alpha_i)$ for $i \in [h]$, $0 \leq j < \mu_i$.
\end{lemma}

\paragraph{Use in the GIJS distinguisher.}
Let $Y = [\mathbf{y}_1 \mid \cdots \mid \mathbf{y}_{n'}] \in
    \mathbb{F}_q^{k' \times n'}$ be the input matrix and fix a hold-out column
$\tau$.  The distinguisher computes the kernel
\[
    K_\tau := \Bigl\{H \in \mathcal{H}_{k',d} :
    J_{\mathbf{y}_j}^{<s}(H) = 0 \text{ for all } j \neq \tau\Bigr\},
\]
where $\mathcal{H}_{k',d}$ is the space of \emph{homogeneous} degree-$d$
polynomials in $k'$ variables, of dimension $\binom{k'+d-1}{d}$.  It outputs
``GRS'' iff every $H \in K_\tau$ vanishes at $\mathbf{y}_\tau$.  Restricting
to homogeneous polynomials kills the unknown GRS column multipliers: for
homogeneous $H$ of degree $d$,
\[
    (\partial^{[\mathbf{a}]} H)(\lambda \mathbf{x})
    = \lambda^{d - |\mathbf{a}|} \cdot (\partial^{[\mathbf{a}]} H)(\mathbf{x})
    \qquad (\lambda \neq 0),
\]
so $J_{\lambda \mathbf{x}}^{<s}(H) = 0 \iff J_{\mathbf{x}}^{<s}(H) = 0$
\cite[Lem.~3.1]{GIJS26}.  The two feasibility conditions are the
\emph{reconstruction} inequality $s(n'-1) > dD$, which forces every
$H \in K_\tau$ to vanish at a clean hold-out via Hermite interpolation along
the hidden degree-$D$ curve, and the \emph{kernel} inequality
$(n'-1)\binom{k'+s-1}{s-1} < \binom{k'+d-1}{d}$, which guarantees
$K_\tau \neq \{0\}$ on a random input so that Schwartz--Zippel applies.


\subsection{The GIJS distinguisher}\label{sec:gijs-dist}


The Ghoshal--Ishai--Jain--Sun~\cite{GIJS26} distinguisher tests whether a given matrix spans
a subcode of a generalized Reed--Solomon code. Given $n'$ columns in
$\mathbb{F}_q^{k'}$, it interpolates: it seeks a nonzero homogeneous degree-$d$
polynomial vanishing to order $s$ at $n'-1$ of the columns (the interpolation
set), then checks whether it also vanishes at the remaining hold-out column.
For a true GRS subcode, the hidden rational-curve structure forces this
vanishing; for random columns, Schwartz--Zippel says it fails with high
probability. The two inequalities below encode feasibility: the first ensures
enough interpolation constraints to force vanishing on the curve, the second
that the interpolation kernel is nonzero.

\begin{theorem}[Ghoshal et al., Theorem~3.2]\label{thm:ghoshal}
    Let $n' \leq q$, $k' \leq K$, $d < q$, $D := K-1$. If
    \begin{align}
        s(n'-1)                   & > d\cdot D \label{eq:ghoshal1}          \\
        (n'-1)\binom{k'+s-1}{s-1} & < \binom{k'+d-1}{d} \label{eq:ghoshal2}
    \end{align}
    then there is a $\mathrm{poly}\!\left(\binom{k'+d-1}{d}, n', \log
        q\right)$-time
    algorithm distinguishing a random $k'$-dimensional subcode of random
    $GRS_{n',K}$
    from uniform random, with gap $\geq 1 - d/q$.
\end{theorem}
The algorithm (Algorithm~1 of~\cite{GIJS26}): hold out an arbitrary column
$\tau$, compute
the kernel $K_\tau$ of Hasse-jet evaluations (order $< s$) at the other $n'-1$
columns in the space of homogeneous degree-$d$ polynomials, and output ``GRS''
iff every $H \in K_\tau$ vanishes at $y_\tau$, ``Random'' otherwise.

\paragraph{Variables.} $k'$ is the dimension of the subcode; $K$ the
hidden ambient GRS dimension ($k' + w$ here); $s$ the jet multiplicity; $d$ the
test-polynomial degree. The first inequality is \emph{reconstruction}; the
second the \emph{kernel} condition.








\section{Phase 1}\label{sec:phase1}
Here we describe the first phase of the attack, which uses random code shortenings and the GIJS distinguisher to 
recover and unmix all the noisy pairs. Given the $k \times (n+w)$ RLCE public matrix, we proceed as follows.
Let $(n', k', s, d, D, q)$ be GIJS parameters
satisfying the kernel inequality
$(n'-1)\binom{k'+s-1}{s-1} < \binom{k'+d-1}{d}$. Here, $D = k'+w$ is the degree
of the hidden rational curve.

\subsection{Choosing the GIJS parameters}\label{sec:param-choice}

Given the RLCE parameters $(n, k, w, q)$ with public-key length $N = n + w$,
the attacker first fixes a target shortened length $n'$. Shortening at
$N - n'$ positions drops the dimension by the same amount, so
\[
    k' = k - (N - n').
\]
After shortening, the hidden ambient GRS code has dimension $K = k' + w$, so the
hidden rational curve has degree
\[
    D = K - 1 = k' + w - 1.
\]
It remains to choose the jet multiplicity $s$ and the test-polynomial degree
$d$. We search over small integer pairs with $1 \le s \le d < q$ satisfying
both GIJS feasibility inequalities:
\begin{enumerate}
    \item \emph{Reconstruction:} $s(n'-1) > d \cdot D$, which forces every
          $H \in K_\tau$ to vanish at a clean hold-out;
    \item \emph{Kernel:} $(n'-1)\binom{k'+s-1}{s-1} < \binom{k'+d-1}{d}$,
          which guarantees the interpolation kernel is nonzero on random input.
\end{enumerate}
Among feasible pairs $(s,d)$, we prefer smaller $d$: the GIJS matrix has
$\binom{k'+d-1}{d}$ columns, so $d$ dominates the per-run cost. Subject to
that constraint, we prefer a larger relative margin on the kernel inequality,
since a
near-tight kernel is numerically fragile.

There is a tension in choosing $n'$ itself. A smaller $k'$ (hence smaller
$n'$) makes each GIJS run cheaper, but GIJS also needs a hard column-count
floor: $n'$ must be large enough that the reconstruction inequality can hold
at all. In practice $n'$ cannot be pushed far below the boundary where
$s(n'-1) \approx dD$.


\subsection{The Core of the Attack}\label{sec:fullattack}
Let $t'$ denote the number
of noisy columns among the $n'-1$ interpolation columns (all columns except the
holdout column) at any given point in the following algorithm.

\paragraph{Dirty column tolerance.}
We define $t'_{\max}$ to be the maximum number of noisy columns such that PIR
persistence (see Section~\ref{sec:gijs}) guarantees acceptance on a clean
hold-out:
\begin{equation}\label{eq:tmax}
    t'_{\max} = \max\{t' : s\cdot((n'-1) - t') > d\cdot D\}.
\end{equation}
A trial is \emph{recorded} --- its flagged set contributes to the
frequency table --- iff $|V| \leq T_{\mathrm{reject}} := t'_{\max}$.

We deliberately obfuscate an inner step, namely the computation of the set of 
holdout columns $\tau$ on which the GIJS subroutine outputs random. 
While this can be trivially computed by separately running the GIJS subroutine 
as a black box on each holdout column, we have a quicker method, described in 
Section~\ref{sec:reuse-attack}.

\medskip

\begin{algorithm}[H]
    \caption{Phase 1}
    \label{alg:main_analysis}
    \KwIn{Original code parameters $n, w, k, k', q$; raw trial count
        $T$; recording threshold $T_{\mathrm{reject}}$}
    \KwOut{All $w$ pairs $(c_1, c_2)$ of noisy columns in the original code}

    $W \leftarrow \emptyset$\;
    \For{$iter \leftarrow 1$ \KwTo $T$}{
        Choose a random subset $\mathcal{L} \subset [1,n+w]$ where
        $\vert{}\mathcal{L}\vert{} = k - k'$\;
        Compute the $k' \times n'$ shortened code by shortening at
        $\mathcal{L}$\;
        $V \leftarrow \emptyset$\;
        Compute $V$ as the set of holdout columns $\tau \in \{1, \dots n'\}$ on which 
        the GIJS subroutine outputs ``random.'' (see Algorithm~\ref{alg:reuse})\\
        $W \leftarrow W \cup \{V\}$\;
        
    }
    Use $W$ to compute all $w$ pairs of noisy columns (see Algorithm~\ref{alg:cluster})\\
    \Return $w$ twin pairs $(c_1, c_2)$
\end{algorithm}

\subsection{Reuse Speedup}\label{sec:reuse-attack}
We can speed up the attack by noticing that the jet matrix for 
each holdout attack has only small differences from the others. Consider a specific shortening. In the language of the GIJS attack 
we have one shortened point set
$y_1,\dots,y_{n'} \in \Fq^{k'}$. For a hold-out $\tau$, the GIJS
decision is computed from the jet matrix $M_\tau$: its rows are
indexed by $(j,\mathbf{a})$ with $j\neq \tau$ and
$\mathbf{a}\in A_{<s,k'}$, its columns by the monomials
$\mathbf{X}^{\mathbf{m}}$ of degree $d$, and its entry is the Hasse
derivative $(\partial^{[\mathbf{a}]}\mathbf{X}^{\mathbf{m}})(y_j)$.
Writing $z$ for the coefficient vector of $H\in V_d$, we have that $(M_\tau
z)_{(j,\mathbf{a})} = (\partial^{[\mathbf{a}]}H)(y_j)$, so
$\ker M_\tau$ is the interpolation kernel $K_\tau$. Let $\mathbf{v}_\tau
\in \Fq^{N_d}$ be the evaluation vector
$\mathbf{v}_\tau = (y_\tau^{\mathbf{m}})_{\mathbf{m}}$ at which the GIJS 
distinguisher is executed.

Now, let $M_{\mathrm{full}}$ be the jet matrix of \emph{all} $n'$ points,
with $E_F := n'J$ rows grouped in $n'$ consecutive blocks of $J$ rows
(block $j$ holds the jets at $y_j$, with $\mathbf{a}=\mathbf{0}$
first). Note that $M_\tau$ is $M_{\mathrm{full}}$ with block $\tau$
deleted, and that $\mathbf{v}_\tau$ is exactly the
$(\tau,\mathbf{0})$-row of $M_{\mathrm{full}}$, since
$(\partial^{[\mathbf{0}]}\mathbf{X}^{\mathbf{m}})(y_\tau) =
y_\tau^{\mathbf{m}}$.

Call a vector $\mathbf{s} = M_{\mathrm{full}} z$ a \emph{syndrome}. As
$z$ ranges over $\Fq^{N_d}$, the syndromes range over the column
space of $M_{\mathrm{full}}$, which is the kernel of any matrix
$\Pperp$ whose rows form a basis of the \emph{left} kernel of
$M_{\mathrm{full}}$ (the row relations $\mathbf{u}^{T} M_{\mathrm{full}} =
0$). Write $h := E_F - \rk(M_{\mathrm{full}})$ for the number of such
basis rows.

\begin{theorem}[Exactness of hold-out reuse]\label{thm:reuse-exact}
For every $\tau\in[n']$, the following are equivalent:
\begin{enumerate}
\item[(i)] the per-hold-out GIJS decision flags $\tau$;
\item[(ii)] there is a syndrome $\mathbf{s} = M_{\mathrm{full}} z$
supported only on block $\tau$ with $s_{(\tau,\mathbf{0})} \neq 0$;
\item[(iii)] $\rk(\Pperp|_\tau) = \rk(\Pperp|_\tau^{-\mathbf{0}})$,
where $\Pperp|_\tau$ is $\Pperp$ restricted to the $J$ columns of
block $\tau$ and $\Pperp|_\tau^{-\mathbf{0}}$ is the same restriction
with the $\mathbf{a}=\mathbf{0}$ column removed. If $\Pperp$ is empty
($h=0$), then \textup{(iii)} is declared to hold for every $\tau$.
\end{enumerate}
Hence all $n'$ decisions of a trial are obtained from one left-kernel
basis $\Pperp$ and two rank computations per hold-out.
\end{theorem}

\begin{proof}
(i)$\iff$(ii): $\tau$ is flagged iff there is $z\in\ker M_\tau$ with
$\langle z,\mathbf{v}_\tau\rangle \neq 0$. For such a $z$, the
syndrome $\mathbf{s} = M_{\mathrm{full}}z$ vanishes outside block
$\tau$ (that is $M_\tau z = 0$) and
$s_{(\tau,\mathbf{0})} = \langle z, \mathbf{v}_\tau\rangle \neq 0$.
Conversely, if $\mathbf{s} = M_{\mathrm{full}} z$ is supported on
block $\tau$ with $s_{(\tau,\mathbf{0})}\neq 0$, then $M_\tau z = 0$
and $\langle z,\mathbf{v}_\tau\rangle = s_{(\tau,\mathbf{0})} \neq 0$.

(ii)$\iff$(iii): the syndromes supported on block $\tau$ are exactly
the vectors $\mathbf{u}\in\Fq^{J}$ with $\Pperp|_\tau \mathbf{u} = 0$,
i.e.\ $\ker(\Pperp|_\tau)$, embedded in block $\tau$: indeed
$\mathbf{s}\in\ker\Pperp$ iff $\mathbf{s} = M_{\mathrm{full}} z$ for
some $z$. Such a pattern with $u_{\mathbf{0}} \neq 0$ exists iff the
first column of $\Pperp|_\tau$ lies in the span of the remaining
columns: if $\mathbf{u}\in\ker(\Pperp|_\tau)$ has
$u_{\mathbf{0}}\neq 0$, then $\mathrm{col}_{\mathbf{0}} =
-\sum_{i\geq 1} (u_i/u_{\mathbf{0}})\,\mathrm{col}_i$; if
$\mathrm{col}_{\mathbf{0}} = \sum_{i\geq 1} w_i\,\mathrm{col}_i$,
then $(1,-w_1,\dots)\in\ker(\Pperp|_\tau)$ has nonzero first entry.
And $\mathrm{col}_{\mathbf{0}}$ lies in the span of the others iff
removing it does not drop the rank, which is (iii). If $h=0$, then
$M_{\mathrm{full}}$ has full row rank, $\ker\Pperp = \Fq^{E_F}$, the
pattern $\mathbf{e}_{\mathbf{0}}$ is achievable on block $\tau$, and
(ii) holds for every $\tau$ directly.
\end{proof}

Algorithm~\ref{alg:reuse} turns Theorem~\ref{thm:reuse-exact} into a three phase
procedure. In the first
phase it assembles the single matrix $M_{\mathrm{full}}$.
Next, it extracts \emph{every} linear relation
among the rows of $M_{\mathrm{full}}$, in one pass. Forward
elimination is run on the augmented matrix
$[M_{\mathrm{full}} \mid I_{E_F}]$. The identity block
is used as a ledger to record combinations
of the original rows formed by forward elimination. Whenever a row of the
left block has been reduced to zero, the ledger row opposite it is
therefore an explicit combination of the rows of
$M_{\mathrm{full}}$ that cancels them.
Collecting all such ledger rows gives the basis $\Pperp$ of the left
kernel. This elimination step has a cost comparable
to one per-hold-out GIJS run. The third
phase, the loop over $\tau$, merely reads $\Pperp$ locally, one block
at a time.

\begin{algorithm}[H]
\caption{Hold-out reuse: all GIJS decisions of one trial}
\label{alg:reuse}
\KwIn{Shortened columns $y_1,\dots,y_{n'}\in\Fq^{k'}$; GIJS
parameters $(s,d)$.}
\KwOut{The set $V\subseteq[n']$ of flagged hold-outs, identical to
$n'$ separate GIJS runs.}
Build $M_{\mathrm{full}}$ ($E_F \times N_d$), rows grouped in blocks
of $J$, $\mathbf{a}=\mathbf{0}$ first in each block\;
Forward-eliminate the augmented matrix $[M_{\mathrm{full}} \mid
I_{E_F}]$; let $\Pperp$ be the rows of the accumulated transform
opposite the zero rows of the echelon form\;
\If{$h = 0$}{\Return $V = [n']$\;}
$F \leftarrow \emptyset$\;
\For{$\tau \leftarrow 1$ \KwTo $n'$}{
    $A \leftarrow \Pperp$ restricted to the $J$ columns of block
    $\tau$\;
    \If{$\rk(A) = \rk(A \text{ without its first column})$}{
        $V \leftarrow V \cup \{\tau\}$\;
    }
}
\Return $V$\;
\end{algorithm}
\smallskip
Restricting $\Pperp$ to the
$J$ columns of block $\tau$ discards everything a relation says
about the other points and keeps only its entries on $\tau$'s own
rows. According to Theorem~\ref{thm:reuse-exact}, $\tau$ is flagged by the GIJS distinguisher
as random if and only if
the first of these $J$ columns is already a linear
combination of the remaining $J-1$. 
As a result, computing the rank of the matrix of the full $J$ columns and comparing it to 
the rank of the last $J-1$ columns is sufficient; the rank drops if the distinguisher outputs GRS
while it stays the same if it outputs random.

The early exit for
$h = 0$ is a degenerate case: with no
row relations at all, every syndrome is achievable and thus every hold-out
is flagged as random.

\subsection{Unmixing Pairs Via Clustering}\label{sec:cluster}
We now show how to take flagged sets of columns from the tria runs and recover
pairs. Here, the flagged sets of columns mean the set $W$
to which sets of columns flagged as ``random" were added in
Algorithm~\ref{alg:main_analysis}
after each good trial.

The idea is simple once the flagged set of a single trial is looked
at the right way. In a trial shortening, a twin pair can be in only three
states. If both members are kept, then both are dirty columns, so
both are flagged with very very high probability.
If exactly one member is kept, and then the survivor
is derandomized to a clean column per the analysis of the Couvrier et al
paper~\cite{CLT18},
so that column is not flagged, and the
other member is absent from the trial altogether. The final case is where both
members
are hit by the shortening, and neither appears. In all three cases
it is with very very high probability the case that either both are flagged or neither
are flageed.

Finally, a column that belongs to no twin pair, (a pure GRS column), is never
flagged at all.

As a result, if both members of a pair are correctly flagged with high enough
probability,
then the list of trials in
which column A was flagged  is the same as the list of columns for which its
twin column B was flagged.
Meanwhile, it is an empty list for a pure column, and for other twin pairs, the
list is non-empty but with
high probability over the random choices of shortenings distinct. Pair
identification thus becomes grouping
columns by their \emph{pattern}.

\subsubsection{The algorithm}

For each original public column $i$, record its pattern
$b_i \in \{0, 1\}^T$, where $b_i(t) = 1$ if $i$ was flagged in
trial $t$ and $0$ otherwise. Then proceed as follows.

\begin{algorithm}[H]
    \caption{Flag-pattern clustering}
    \label{alg:cluster}
    \KwIn{Flagged sets $V_{1}, \dots, V_{T}$ of $T$ trials,
        mapped back
        to original public-column indices}
    \KwOut{Sets of twin pairs and pure columns; anomalies, if any}
    For each column $i$, form its pattern $b_i$, with $b_i(t) = 1$
    iff $i \in V_t$\;
    Group columns by \emph{exact equality} of patterns\;
    A group of size $2$ with a nonzero pattern is a twin pair\;
    The group with the all-zero pattern is the pure columns\;
    A group of size $1$ with a nonzero pattern is an orphan\;
    A group of size $3$ or more is oversized, and is reported, not
    paired\;
    Compare the size of the all-zero group with the expected
    number $n - w$ of pure columns, and report any excess\;
\end{algorithm}



\subsubsection{Trials for Success}\label{sec:cluster-planning}

We now compute how many trials (assuming that pairs are with overwhelming probability correctly flagged together) 
are necessary to identify and unmix pairs with high probability. Fix a pair. 
The probability both columns of the pair appear in the shortening is \[
    p \;=\; \frac{(n')(n'-1)}{(n+w)(n+w-1)},
\]
the probability that a given pair survives a random shortening to
$n'$ out of $N$ columns. Two distinct pairs separate in a recorded trial
except with probability at
most $p^2 + (1 - p)^2$, when either both pairs appear in the shortening, or neither pair appears in the shortening.


Note that we heuristically treat the pairs as independent here. A union bound over the $\binom{w}{2}$ pairs of
pairs then says that $T$ trials separate every pair from every
other pair with failure probability at most $\epsilon$ as soon as
\[
    \binom{w}{2} \bigl(p^2 + (1 - p)^2\bigr)^T \;\leq\; \epsilon,
\]

and we take $T_{\mathrm{sep}}$ to be the smallest such $T$. A pair
is invisible, and so hides in the pure set, if it is never unhit
in $T$ trials, which happens with probability $(1 - p)^T$ per pair;
a union bound over the $w$ pairs gives the visibility requirement
\[
    w (1 - p)^T \;\leq\; 0.01,
\]
with $T_{\mathrm{vis}}$ the smallest such $T$. Clustering runs $
    T \;=\; \max(T_{\mathrm{sep}}, T_{\mathrm{vis}})
$.

\subsubsection{Empirical Validation}\label{sec:cluster-validation}

Clustering was validated at several points, where the truth
was known from the key and every trial could be checked against
it. 

These parameters include the actual ID-0 parameters. 

\begin{table}[htbp]
    \centering
    \footnotesize
    \begin{tabular}{lcccc}
        \toprule
        Scheme Parameters ($q, n, k, w$) & Attack Parameters ($n',(s,d)$)  & Trials & Flagging exact & Pairs recovered \\
        \midrule
        $(256, 198, 165, 33)$, $(90, 2, 3)$           & $23/23$        & $33/33$    \\
        $(256, 119, 79, 31)$, $(101, 2, 3)$  & $25/25$        & $31/31$   \\
        $(1024, 630, 470, 160)$, $(366, 2, 3)$ & $41/41$ & $160/160$ \\ 
        \bottomrule
    \end{tabular}
    \caption{Cluster validation at reduced parameters. ``Flagging
        exact'' means the flagged set equalled the unhit pairs' columns
        in that many trials.}
    \label{tab:cluster-validation}
\end{table}


\paragraph{ID-0, end to end.} The shipped ID-0 set has since been
run end to end with cluster Phase~1, on a fresh key generated with
the NIST DRBG key generation of the implementation, at the adopted
point $(n', k', (s, d)) = (366, 46, (2, 3))$, with $41$ trials on
$12$ threads. The run succeeded completely: cluster aggregation
over the $41$ trials returned $160$ pairs, $470$ pure columns,
$0$ orphans, and $0$ oversized groups; Phase~2 unmixed $160/160$
pairs, Sidelnikov--Shestakov recovered the correct GRS code, and
verification passed with $160/160$ non-degenerate pairs correct
and $0$ false positives. The run took a little over $3$ hours of
wall time across those $12$ threads to break ID-0. That is the
Phase~1 trial time dominating the run: $41$ trials in batches of
$12$, each trial a single reuse elimination, with aggregation,
Phase~2 (SS bootstrap $0.25$\,s), and verification taking minutes
by comparison. ID-1 and ID-2 have not been run end to end with
cluster Phase~1; their figures below remain model estimates.




\section{The GIJS Distinguisher In the Our Context}\label{sec:gijs}
In this section we analyze the performance of the GIJS
distinguisher~\cite{GIJS26} in our attack algorithms.

We now show that under plausible and empirically validated heuristics, the GIJS
distinguisher algorithms succeeds in distinguishing
clean columns (shortened from a pure GRS column) from dirty columns (one of two
elements of a twin pair, both of which appear in the shortened code) under
certain
conditions even with noisy RLCE columns.

The shortened clean columns lie on a projected
rational curve
\[
    \varphi : \mathbf{P}^1 \longrightarrow \mathbf{P}^{k'-1},
\]
of degree $D=k'+w-1$, given by homogeneous polynomials of degree $D$ in the
decoding
parameter $t$. Below, we work over the algebraic closure
$\overline{\mathbf{F}}_q$ when counting roots.

Each original column has a distinct evaluation parameter
$\alpha \in \mathbf{F}_q$, and this distinct parameter is unchanged by
shortening.

Fix a holdout column and let $I$ be the set of interpolation
columns, $|I| = n'-1$. Partition $I$ as follows:
\begin{itemize}
    \item \emph{Clean points:} kept pure columns and singleton-hit
          survivors. Each lies on the curve $\varphi(\mathbf{P}^1)$,
          at the parameter of its original column. Let $\mathcal{C}$
          denote this set, and $C = |\mathcal{C}|$.
    \item \emph{Dirty points:} kept columns belonging to a pair both
          of whose members are kept. Such a pair is called
          \emph{both-kept}. Let $\mathcal{U}$ denote the set of
          non-degenerate both-kept pairs, and $U = |\mathcal{U}|$.
          Non-degenerate means that neither mixing coefficient is zero,
          so neither member of the pair is itself a curve point; in the
          RLCE specification the mixing coefficients are chosen to always
          be non-degenerate.
\end{itemize}
Thus, counting kept columns, $n' = C + 2U$ plus any degenerate
both-kept pairs, which we exclude from $\mathcal{U}$ and treat
separately; in the ID-0 run below there are none.

Let $H$ be a homogeneous polynomial of degree $d$ in $k'$
variables. We say $H$ \emph{vanishes to order $s$} at a point
$x$ if all partial derivatives of $H$ of order $< s$ vanish at $x$,
i.e.\ the jet of $H$ at $x$ of order $s-1$ is zero. The GIJS
kernel, for a fixed interpolation set, is the space of such $H$
vanishing to order $s$ at every point of $I$. A clean holdout
$\tau$ is accepted, i.e.\ not flagged, if every $H$ in the kernel
vanishes at $\tau$; it is flagged if some $H$ in the kernel does
not vanish at $\tau$.

\subsection{PIR persistence: rigorous acceptance with
    noise}\label{sec:pir-noise}

First, we show that the pairs of dirty columns (which are the only type of
non-clean columns present in the shortened code)
are actually useful to the GIJS attack. As we will show, the reason is
precisely due to the linear relation existing between pairs that results in
a pure GRS column when inverting the $A$ map from public key generation.

\begin{theorem}[PIR persistence]
    Let $H$ be a homogeneous polynomial of degree $d$ vanishing to
    order $s$ at every interpolation point, and assume
    \begin{equation}
        2s > d.
        \label{eq:line}
    \end{equation}
    Then $H$ vanishes identically on the line through the two dirty
    columns of every non-degenerate both-kept pair, and in particular
    $H$ vanishes at the underlying GRS point of that pair.
    Consequently, if
    \begin{equation}
        sC + U > dD,
        \label{eq:condition}
    \end{equation}
    then $H$ vanishes identically on the curve $\varphi(\mathbf{P}^1)$.

    As a result, it vanishes at every clean holdout column, meaning that
    running
    the
    GIJS distinguisher with a clean holdout column outputs "GRS"
    deterministically
    if the above conditions are fulfilled.
\end{theorem}

\begin{proof}
    Fix a non-degenerate both-kept pair with underlying GRS point
    $P$ and random point $R$. After shortening, its two columns are
    \[
        Q_1 = aP + cR, \qquad Q_2 = bP + dR,
    \]
    in projective coordinates, with $ad - bc \ne 0$ and $c, d \ne 0$
    by non-degeneracy. Hence $Q_1 \ne Q_2$ in $\mathbf{P}^{k'-1}$,
    and $Q_1, Q_2, P$ are collinear on the line $L$ spanned by $P$
    and $R$; $L$ is not contained in the curve, since $R$ is not a
    curve point.

    Restrict $H$ to $L$. The restriction is a binary form of degree
    $d$ on $L$, vanishing to order $s$ at the two distinct points
    $Q_1$ and $Q_2$. A binary form of degree $d$ with two distinct
    roots of multiplicity $s$ each is identically zero as soon as
    $2s > d$, by degree counting. This is~\eqref{eq:line}. Hence
    $H$ vanishes identically on $L$, and in particular $H(P) = 0$.

    Now consider the pullback $h(t) = H(\varphi(t))$. This is a binary
    form of degree $dD$, unless it is identically zero, which is the
    conclusion we want. We count its roots over
    $\overline{\mathbf{F}}_q$, with multiplicity.

    First, let $t_0$ be the parameter of a clean interpolation point
    $X = \varphi(t_0)$. Since $H$ vanishes to order $s$ at $X$ in
    ambient space, its Taylor expansion at $X$ has no terms of degree
    $< s$. Composing with $\varphi(t) - X$, which has zero constant
    term at $t_0$, gives a power series in $t - t_0$ with no terms of
    degree $< s$; if $\varphi$ is ramified at $t_0$, the order is only
    larger. Hence $h$ vanishes to order at least $s$ at $t_0$.

    Second, let $t_i$ be the parameter of the underlying GRS point
    $P_i$ of a non-degenerate both-kept pair. By the line argument
    above, $H(P_i) = 0$, so $h(t_i) = 0$, i.e.\ $h$ has a root at
    $t_i$ of order at least $1$. The parameters $t_i$ are distinct
    across pairs, since the original evaluation parameters are
    distinct, and they are distinct from the clean parameters, since
    a both-kept pair's parameter is not the parameter of any clean
    interpolation column. Hence the roots counted in the first and
    second steps are distinct, except that a clean root already
    carries multiplicity $s$.

    Therefore $h$ has total root multiplicity at least $sC + U$.
    If \eqref{eq:condition} holds, this exceeds the degree $dD$ of
    $h$, so $h$ cannot be a nonzero binary form. Hence $h \equiv 0$,
    i.e.\ $H$ vanishes at every point of the curve
    $\varphi(\mathbf{P}^1)$. A clean holdout column is a curve
    point, so every kernel polynomial $H$ vanishes at it, and the
    GIJS test accepts it.
\end{proof}

\subsection{Random Holdout Columns}\label{sec:rand-holdout}
We now show that with high probability (and heuristically with overwhelming
probability) that a random and independent column
will be flagged by the GIJS distinguisher, meaning that at least one kernel
polynomial $H$ will not vanish at it.

\begin{theorem}[Independent uniform random holdout]
    Fix the interpolation set $I$, and let $K$ be its GIJS kernel
    as above. Assume $K \ne \{0\}$. Let $\tau$ be uniform in
    $\mathbf{F}_q^{k'}$, independent of $I$ and hence independent
    of $K$. Then
    \begin{equation}
        \Pr[\tau \text{ is flagged}] \ge 1 - \frac{d}{q}.
        \label{eq:bound}
    \end{equation}
\end{theorem}

\begin{proof}
    Since $K \ne \{0\}$, pick any nonzero $H_0 \in K$. The flagging
    event contains the event $\{H_0(\tau) \ne 0\}$, because if
    $H_0$ does not vanish at $\tau$, the GIJS test has a witness
    and flags $\tau$. Hence
    \[
        \Pr[\tau \text{ is flagged}]
        \ge \Pr[H_0(\tau) \ne 0]
        = 1 - \Pr[H_0(\tau) = 0].
    \]
    It remains to bound $\Pr[H_0(\tau) = 0]$. This is the
    Schwartz--Zippel bound: a nonzero polynomial of total degree
    $d$ in $k'$ variables over $\mathbf{F}_q$ vanishes at at most
    a $d/q$ fraction of points in $\mathbf{F}_q^{k'}$. Indeed, the
    bound follows by induction on $k'$, writing $H_0$ as a
    polynomial in one variable with coefficients in the remaining
    variables; a univariate polynomial of degree at most $d$ has at
    most $d$ roots unless it is identically zero. Applying it to
    $H_0$ gives $\Pr[H_0(\tau) = 0] \le d/q$, and
    \eqref{eq:bound} follows.
\end{proof}

\subsubsection{Non-zero Kernel Conditions}\label{sec:nonzero-kernel}
It remains to discuss conditions on which the kernel is non-zero. 

Let $J=\binom{k'+s-1}{s-1}$ be the number of jet conditions, and 
$N_d=\binom{k'+d-1}{d}$ the number of monomials of degree $d$ in $k'$ variables.

\begin{corollary}[Sufficient condition for a nonzero kernel]
    If
    \begin{equation}
        (n'-1)J < N_d,
        \label{eq:kernel}
    \end{equation}
    then $K \ne \{0\}$ for every interpolation set, and hence
    \eqref{eq:bound} holds for every trial.
\end{corollary}

\begin{proof}
    $K$ is the kernel of an $(n'-1)J \times N_d$ matrix. Its rank
    is at most the number of rows, $(n'-1)J$. By rank--nullity,
    \[
        \dim K \ge N_d - (n'-1)J > 0
    \]
    under \eqref{eq:kernel}, so $K \ne \{0\}$. Theorem~3 then
    applies.
\end{proof}


\subsubsection{Heuristic Bounds on Dirty Columns}\label{sec:heuristic}

As our experiments (see below) show, we have not had \emph{any} dirty columns 
fail to flag (i.e. the GIJS distinguisher outputs "GRS" when they are a holdout column). 

Under a reasonable heuristic and using the $N_d - (n'-1)J$ lower bound for the kernel dimension, it seems 
reasonable to assume that the probability of failure in practice is more on the order of $(d/q)^{\dim K}$, 
meaning the probability is essentially independent across kernel polynomials. We stress
that this is a pure heuristic, but that it has been experimentally shown to be reasonable on 
parameter sets up to the lowest set given in the RLCE specification. 



% \input{puncture}



% \subsection{Unmixing Pairs Via Clustering}\label{sec:cluster}
We now show how to take flagged sets of columns from the tria runs and recover
pairs. Here, the flagged sets of columns mean the set $W$
to which sets of columns flagged as ``random" were added in
Algorithm~\ref{alg:main_analysis}
after each good trial.

The idea is simple once the flagged set of a single trial is looked
at the right way. In a trial shortening, a twin pair can be in only three
states. If both members are kept, then both are dirty columns, so
both are flagged with very very high probability.
If exactly one member is kept, and then the survivor
is derandomized to a clean column per the analysis of the Couvrier et al
paper~\cite{CLT18},
so that column is not flagged, and the
other member is absent from the trial altogether. The final case is where both
members
are hit by the shortening, and neither appears. In all three cases
it is with very very high probability the case that either both are flagged or neither
are flageed.

Finally, a column that belongs to no twin pair, (a pure GRS column), is never
flagged at all.

As a result, if both members of a pair are correctly flagged with high enough
probability,
then the list of trials in
which column A was flagged  is the same as the list of columns for which its
twin column B was flagged.
Meanwhile, it is an empty list for a pure column, and for other twin pairs, the
list is non-empty but with
high probability over the random choices of shortenings distinct. Pair
identification thus becomes grouping
columns by their \emph{pattern}.

\subsubsection{The algorithm}

For each original public column $i$, record its pattern
$b_i \in \{0, 1\}^T$, where $b_i(t) = 1$ if $i$ was flagged in
trial $t$ and $0$ otherwise. Then proceed as follows.

\begin{algorithm}[H]
    \caption{Flag-pattern clustering}
    \label{alg:cluster}
    \KwIn{Flagged sets $V_{1}, \dots, V_{T}$ of $T$ trials,
        mapped back
        to original public-column indices}
    \KwOut{Sets of twin pairs and pure columns; anomalies, if any}
    For each column $i$, form its pattern $b_i$, with $b_i(t) = 1$
    iff $i \in V_t$\;
    Group columns by \emph{exact equality} of patterns\;
    A group of size $2$ with a nonzero pattern is a twin pair\;
    The group with the all-zero pattern is the pure columns\;
    A group of size $1$ with a nonzero pattern is an orphan\;
    A group of size $3$ or more is oversized, and is reported, not
    paired\;
    Compare the size of the all-zero group with the expected
    number $n - w$ of pure columns, and report any excess\;
\end{algorithm}



\subsubsection{Trials for Success}\label{sec:cluster-planning}

We now compute how many trials (assuming that pairs are with overwhelming probability correctly flagged together) 
are necessary to identify and unmix pairs with high probability. Fix a pair. 
The probability both columns of the pair appear in the shortening is \[
    p \;=\; \frac{(n')(n'-1)}{(n+w)(n+w-1)},
\]
the probability that a given pair survives a random shortening to
$n'$ out of $N$ columns. Two distinct pairs separate in a recorded trial
except with probability at
most $p^2 + (1 - p)^2$, when either both pairs appear in the shortening, or neither pair appears in the shortening.


Note that we heuristically treat the pairs as independent here. A union bound over the $\binom{w}{2}$ pairs of
pairs then says that $T$ trials separate every pair from every
other pair with failure probability at most $\epsilon$ as soon as
\[
    \binom{w}{2} \bigl(p^2 + (1 - p)^2\bigr)^T \;\leq\; \epsilon,
\]

and we take $T_{\mathrm{sep}}$ to be the smallest such $T$. A pair
is invisible, and so hides in the pure set, if it is never unhit
in $T$ trials, which happens with probability $(1 - p)^T$ per pair;
a union bound over the $w$ pairs gives the visibility requirement
\[
    w (1 - p)^T \;\leq\; 0.01,
\]
with $T_{\mathrm{vis}}$ the smallest such $T$. Clustering runs $
    T \;=\; \max(T_{\mathrm{sep}}, T_{\mathrm{vis}})
$.

\subsubsection{Empirical Validation}\label{sec:cluster-validation}

Clustering was validated at several points, where the truth
was known from the key and every trial could be checked against
it. 

These parameters include the actual ID-0 parameters. 

\begin{table}[htbp]
    \centering
    \footnotesize
    \begin{tabular}{lcccc}
        \toprule
        Scheme Parameters ($q, n, k, w$) & Attack Parameters ($n',(s,d)$)  & Trials & Flagging exact & Pairs recovered \\
        \midrule
        $(256, 198, 165, 33)$, $(90, 2, 3)$           & $23/23$        & $33/33$    \\
        $(256, 119, 79, 31)$, $(101, 2, 3)$  & $25/25$        & $31/31$   \\
        $(1024, 630, 470, 160)$, $(366, 2, 3)$ & $41/41$ & $160/160$ \\ 
        \bottomrule
    \end{tabular}
    \caption{Cluster validation at reduced parameters. ``Flagging
        exact'' means the flagged set equalled the unhit pairs' columns
        in that many trials.}
    \label{tab:cluster-validation}
\end{table}


\paragraph{ID-0, end to end.} The shipped ID-0 set has since been
run end to end with cluster Phase~1, on a fresh key generated with
the NIST DRBG key generation of the implementation, at the adopted
point $(n', k', (s, d)) = (366, 46, (2, 3))$, with $41$ trials on
$12$ threads. The run succeeded completely: cluster aggregation
over the $41$ trials returned $160$ pairs, $470$ pure columns,
$0$ orphans, and $0$ oversized groups; Phase~2 unmixed $160/160$
pairs, Sidelnikov--Shestakov recovered the correct GRS code, and
verification passed with $160/160$ non-degenerate pairs correct
and $0$ false positives. The run took a little over $3$ hours of
wall time across those $12$ threads to break ID-0. That is the
Phase~1 trial time dominating the run: $41$ trials in batches of
$12$, each trial a single reuse elimination, with aggregation,
Phase~2 (SS bootstrap $0.25$\,s), and verification taking minutes
by comparison. ID-1 and ID-2 have not been run end to end with
cluster Phase~1; their figures below remain model estimates.




\section{Phase 2: unmixing three pairs}\label{sec:fastp2}

Phase~1 hands us two things: the twin
pairings, and the pure columns. However, this is insufficient for Sidelnikov--Shestakov~\cite{sidelnikov1992insecurity}, 
as that requires $k+3$ pure columns, and at the $w=n-k$ parameters, Phase~1 finds exactly $k$
pure columns. 

A trivial method of finding the pairs would be to compute linear combination $c_{1,i}+a_i*c_{2,i}$ for $a \neq 0$
 for each of 3 
pairs $(c_{1,i}, c_{2,i})$ for which we wish to recover pure columns. Exactly 1 of the $(q-1)^3$ possibilities of 
choices of 3 $a_i$ will 
successfully unmix all the 3 columns, enabling the Sidelnikov--Shestakov attack to succeed.

Instead, we use some clever ideas to \emph{solve} for the triple.
It costs only about $q^2$ cheap steps.

\subsection{Putting the pairs in systematic form}
We begin with the matrix \[G = [ P \mid c_{1,1} \mid c_{1,2} \mid c_{2,1} \mid c_{2,2} \mid c_{3,1} \mid c_{3,2}]\] where
 $P$ is the $k \times k$
matrix of the pure columns. After
a left multiply of $P^{-1}$, the generator looks like
$G=[I_k \mid \mathbf{g}_1\ \mathbf{g}_2\ \mathbf{g}_3]$,
where for pair $i$ with public columns $c_{i,1}, c_{i,2}$ we set
$\mathbf{a}_i = P^{-1} c_{i,1}$,
$\mathbf{b}_i = P^{-1} c_{i,2}$, and
\[
    \mathbf{g}_i(\lambda_i) \;=\; \mathbf{a}_i + \lambda_i \mathbf{b}_i,
    \qquad \lambda_i \in \mathbb{P}^1(\Fq),
\]
with $\mathbf{g}_i(\infty) := \mathbf{b}_i$. 

\subsection{The Test}

The dual $D(\lambda)$ is generated by the rows of
\[
    H \;=\; \bigl[\, -\mathbf{g}_1^{T}\ ;\ -\mathbf{g}_2^{T}\ ;\
    -\mathbf{g}_3^{T} \;\big|\; I_3 \,\bigr]
\]
 For $j \in [k]$
write $\gamma_j = (\mathbf{g}_1[j], \mathbf{g}_2[j],
\mathbf{g}_3[j]) \in \Fq^3$ for the $j$-th column of the left
block, and
\[
    \mathbf{m}_j \;=\; \bigl(\gamma_{j1}\gamma_{j2},\
    \gamma_{j1}\gamma_{j3},\ \gamma_{j2}\gamma_{j3}\bigr) \in \Fq^3,
    \qquad
    M^{*}(\lambda) := \bigl[\mathbf{m}_1 \mid \cdots \mid
    \mathbf{m}_k\bigr] \in \Fq^{3\times k}.
\]
In words, $\mathbf{m}_j$ is just the three pairwise products of
the entries of $\gamma_j$. The test is about the rank of this
$3 \times k$ matrix.

\begin{lemma}[Dual-square characterization]\label{lem:fastp2-dual}
Assume $k \geq 2$. Then
\[
    \dim D(\lambda)^{2} \;=\; 3 + \rk M^{*}(\lambda),
\]
and $C(\lambda)$ is a GRS code with distinct support if and only if
\begin{enumerate}
\item[(i)] $\rk M^{*}(\lambda) = 2$;
\item[(ii)] the (then unique up to scale) vector $\mathbf{c} \in
\Fq^3$ with $\mathbf{c}\, M^{*}(\lambda) = 0$ satisfies $c_{12}
c_{13} c_{23} \neq 0$;
\item[(iii)] the $k+3$ columns of $H$ are pairwise
non-proportional.
\end{enumerate}
\end{lemma}

\begin{proof}
Pairwise products of the three generators of $D(\lambda)$ span
$D(\lambda)^2$. The evaluation map sending a quadric
$Q \in \mathrm{Sym}^2(\Fq^3)$, a space of dimension $6$, to the
vector of its values on the columns of $H$ has as its kernel
exactly the quadrics vanishing at all $k+3$ columns of $H$,
viewed as points of $\mathbb{P}^2$. Vanishing at
$\mathbf{e}_1, \mathbf{e}_2, \mathbf{e}_3$ forces the squared
coefficients of $Q$ to be zero, so every quadric in the kernel
looks like $Q_{\mathbf{c}} = c_{12}X_1X_2 + c_{13}X_1X_3 +
c_{23}X_2X_3$, and $Q_{\mathbf{c}}(\gamma_j) =
\mathbf{c}\cdot\mathbf{m}_j$. Hence the kernel has dimension
$3 - \rk M^{*}$, and $\dim D(\lambda)^2 = 6 - (3 - \rk M^{*})
= 3 + \rk M^{*}$, in any characteristic.

If $C(\lambda)$ is GRS with distinct support, so is its dual
$D(\lambda)$~\cite{MS77}, a $\GRS_{k+3,3}$ whose square is the
evaluation code of degree at most $4$ polynomials, of dimension
$\min(5, k+3) = 5$; hence $\rk M^{*} = 2$. The columns of $H$
are the support points of $D(\lambda)$, hence pairwise
non-proportional, and they all lie on a smooth conic. That conic
is a quadric through $\mathbf{e}_1,\mathbf{e}_2,\mathbf{e}_3$,
so it is $Q_{\mathbf{c}}$ for the unique relation
$\mathbf{c}$, and smoothness is exactly (ii): if, say,
$c_{12} = 0$ then $Q_{\mathbf{c}} = X_3(c_{13}X_1 + c_{23}X_2)$
is a union of two lines, whereas if all $c_{ij} \neq 0$ the
gradient system in characteristic $2$ has the unique candidate
singular point, the nucleus $\nu = (c_{23}, c_{13}, c_{12})$,
and $Q_{\mathbf{c}}(\nu) = c_{12}c_{13}c_{23} \neq 0$, so $\nu$
is not on the conic and the conic is smooth.  While we do not need 
it for the purposes of this paper, the same conclusion
in odd characteristic follows from the determinant
$c_{12}c_{13}c_{23}/4 \neq 0$.

Conversely, given (i)--(iii), all $k+3$ columns of $H$ lie on
$Q_{\mathbf{c}}$, which is a smooth conic by (ii), and a
$3$-dimensional code whose generator columns are pairwise
distinct points of a smooth conic is a $\GRS_{3}$ code, by
parametrizing the conic by $\mathbb{P}^1$. So $D(\lambda)$ is
GRS and $C(\lambda) = D(\lambda)^{\perp}$ is GRS with distinct
support by (iii).
\end{proof}

\subsection{Solving for the triple, instead of trying it}

Now, we present our root-finding idea. Fix a triple $J = \{j_1,j_2,j_3\} \subset [k]$ of coordinate positions (for simplicity we may take $(0,1,2)$), and let
\[
    \Delta_J(\lambda) \;=\; \det\bigl[\, R_{ab}(j_t) \,\bigr]_{
    (a,b)\,\in\,\{(1,2),(1,3),(2,3)\},\ t\in\{1,2,3\}}
\]
be the corresponding $3\times 3$ minor of $M^{*}(\lambda)$,
where $R_{ab}(j) = \mathbf{g}_a[j]\,\mathbf{g}_b[j]$. Each
$R_{ab}$ is a product of two affine forms, so every term of the
determinant has degree at most $2$ in each $\lambda_i$:
$\Delta_J$ has multidegree at most $(2,2,2)$. By the lemma,
$\Delta_J$ vanishes at the true triple $\lambda^{*}$, and
indeed every $3\times 3$ minor of $M^{*}(\lambda^{*})$ does,
since $M^{*}(\lambda^{*})$ has rank $2$.

Fix $(\lambda_1,\lambda_2)$ and call that a \emph{cell}. The
row $(R_{12}(j))_{j\in J}$ does not involve $\lambda_3$; call
it $\mathbf{e}$. The other two rows are affine in $\lambda_3$:
$(R_{13}(j))_{j\in J} = \mathbf{p} + \lambda_3 \mathbf{q}$ and
$(R_{23}(j))_{j\in J} = \mathbf{r} + \lambda_3 \mathbf{s}$.
Multilinearity of the determinant in the last two rows gives
\[
    \Delta_J \;=\; A\,\lambda_3^{2} + B\,\lambda_3 + C,
\]
with $A, B, C$ each a $3\times 3$ determinant in
$\mathbf{e},\mathbf{p},\mathbf{q},\mathbf{r},\mathbf{s}$.
So at each cell, the admissible $\lambda_3$ are just the roots
of one explicit quadratic, costing about $150$ field operations
per cell, independent of $k$. Solving that quadratic over
$\mathbb{F}_{2^m}$ is closed-form: if $A = 0$ it is linear (or
constant); if $A \neq 0$ and $B = 0$ the unique root is
$\lambda = (C/A)^{2^{m-1}}$; if $A,B \neq 0$, substituting
$\lambda = (B/A)\mu$ gives $\mu^2 + \mu = AC/B^2$, solvable
iff $\operatorname{Tr}(AC/B^2) = 0$, with roots $\mu, \mu+1$.
Heuristically, we have observed that the number of roots over all $q^2$ cells is about $q^2$ in
expectation (with $68{,}406$ observed in a test at $q=256$).

The rest is a cascade, and each stage is cheaper than the last
in the only sense that matters, the number of survivors.
\textbf{Stage~1} computes the roots at every cell.
\textbf{Stage~2} evaluates a second minor $\Delta_{J'}$, for a
disjoint triple $J'$, at each root (about $60$ operations),
leaving about $q$ survivors. \textbf{Stage~3} applies the exact
test of Lemma~\ref{lem:fastp2-dual} to each survivor ($O(k)$
operations each).~\textbf{Stage~4} rebuilds each surviving
candidate in original coordinates and runs SS, accepting only
a triple that also passes the attack's deterministic full-key
verification. The remaining $w - 3$ pairs are then unmixed
against the recovered curve, and the key completed, exactly as
in the original Phase~2 (including the CLT
post-processing~\cite{CLT18}); none of that is changed.

\begin{algorithm}[H]
\caption{Root-finding unmixing (one orientation)}
\label{alg:fastp2}
\KwIn{The $k$ pure columns; three pairs $(c_{i,1},c_{i,2})$; disjoint
triples $J, J' \subset [k]$.}
\KwOut{A verified unmixing triple, or failure.}
$P \leftarrow$ pure-column matrix; compute $P^{-1}$ and
$\mathbf{a}_i, \mathbf{b}_i$ for $i=1,2,3$\;
\For{$(\lambda_1,\lambda_2) \in \Fq^2$}{
    Form $\mathbf{e},\mathbf{p},\mathbf{q},\mathbf{r},\mathbf{s}$ on
    $J$; compute $A,B,C$\;
    \eIf{$A = B = C = 0$}{$\mathcal{R} \leftarrow \{\lambda_3 \in
        \Fq : \Delta_{J'}(\lambda_1,\lambda_2,\lambda_3) =
        0\}$ \tcp*{degenerate cell: scan}}
    {$\mathcal{R} \leftarrow \{\lambda_3 : \Delta_{J'}(\lambda_1,
        \lambda_2, \lambda_3) = 0,\ \lambda_3 \text{ a root of }
        A\lambda_3^2{+}B\lambda_3{+}C\}$\;}
    \For{$\lambda_3 \in \mathcal{R}$}{
        \If{$(\lambda_1,\lambda_2,\lambda_3)$ passes
            Lemma~\ref{lem:fastp2-dual} and SS $+$ verification on
            original coordinates}{\Return
            $(\lambda_1,\lambda_2,\lambda_3)$\;}
    }
}
\Return failure
\end{algorithm}


\section{Experimental Costs for Key Recovery}
\label{sec:total-costs}

Experiments were run on a single-socket, single-NUMA-node x86\_64 machine with
an Intel Core Ultra 7 265H processor (family 6, model 197, stepping 2): 16
cores and 16 hardware threads (one thread per core), with a maximum frequency
of 5.3\,GHz, 30\,GiB of RAM and 8\,GiB of swap, a 24\,MiB shared L3 cache,
28\,MiB of L2 cache in total, and 480\,KiB of L1 data cache and 768\,KiB of L1
instruction cache in total.

All experiments below were successful in key recovery. Code will be available on Github shortly if the 
author gets around to it. 

\subsection{Parameter Sets Tested}\label{sec:param-sets}

In addition to the ID-0 (128-bit) parameter set from the RLCE specification, we
also tested several smaller parameter sets, and present them and name them here
to avoid cluttering the timing table below; all sets below have $w=n-k$. 


\begin{table}[!ht]
\caption{Parameter Sets Tested}
\begin{center}
\begin{tabular}{llcccc}
\toprule
Name        & q & n & k & $n'$\\
\midrule
256  & 256 & 198 & 165 & 90\\
512  & 512 & 500 & 375 & 293\\
\bottomrule
\end{tabular}
\end{center}
\end{table}

\subsection{Timing tables}\label{sec:timing}

Below, phase 1 refers to the total time taken summed across threads.
\begin{table}[!ht]
\caption{Timing Results}
\begin{center}
    \footnotesize
    \begin{tabular}{llrrrrr}
        \toprule
        Set                           & Trials             & Threads & Phase~1 Time & Phase~2 Time \\
        \midrule
        256 & 50 & 10 & 288.27 (s) & 0.02 (s)\\
        512 & 44 & 12 & 37278.15 (s) & 0.09 (s) \\
        ID-0    & 41             & 12   & 109673.21 (s)        & 0.23 (s)     \\
        \bottomrule
    \end{tabular}
\end{center}
\end{table}



% \section{The $w=n$ Defense: ISD-Secure Parameters}\label{sec:wn-params}

% Updated: the pairing step below is now clustering
% (\S\ref{sec:cluster}) followed by Phase~2 root-finding
% (\S\ref{sec:fastp2}), not the pair-frequency counting and
% per-pair GIJS disentangling of the earlier version of this
% section. The ISD table itself is unchanged; the cost table
% at the end of this section is the old model, kept only for
% comparison, and its replacement belongs in
% \S\ref{sec:total-costs}, to be filled in from new runs.

With $w=n$ we have $N=2n$ and error weight $t=\lfloor(n-k)/2\rfloor$ (the GRS
decoding radius, unchanged --- decryption unmixes all $n$ pairs and decodes as
before). The ISD work factor drops sharply versus $w=n-k$: $t$ is fixed by
$(n-k)/2$ while $N$ grows to $2n$, so the error density $t/N$ falls.

We searched rates $R=k/n\in[0.40,0.85]$ for minimal $n$ with Stern ISD
at the target level. The Stern work factors are bit-operation counts in
the standard model of Stern's algorithm~\cite{Stern89}: expected
number of iterations times per-iteration cost, minimized over the
zero-window parameters $(p, \ell)$. The optimum sits at $R\approx 0.52$ --- well below the
spec rates ($\approx 0.75$).
\begin{itemize}
    \item \textbf{raw}: minimal $(n,k)$ with Stern $\geq$ target;
    \item \textbf{cons}: Stern $\geq$ target$+15$, to account for improvements such as that of Becker et al~\cite{DBLP:conf/eurocrypt/BeckerJMM12}. 
\end{itemize}
The field is the smallest $\mathrm{GF}(2^m)$ with $n\leq 2^m$.

\begin{center}
    \begin{tabular}{lrrrrrrrr}
        \toprule
                 & $n$  & $k$  & $t$ & $N$  & field    & Stern & pk (KB) \\
        \midrule
        128-raw  & 1080 & 562  & 259 & 2160 & $2^{11}$ & 128.1 & 1630    \\
        128-cons & 1218 & 632  & 293 & 2436 & $2^{11}$ & 143.1 & 2067    \\
        192-raw  & 1665 & 865  & 400 & 3330 & $2^{11}$ & 192.1 & 3868    \\
        192-cons & 1801 & 937  & 432 & 3602 & $2^{11}$ & 207.1 & 4532    \\
        256-raw  & 2246 & 1168 & 539 & 4492 & $2^{12}$ & 256.0 & 7686    \\
        256-cons & 2382 & 1238 & 572 & 4764 & $2^{12}$ & 271.0 & 8639    \\
        \bottomrule
    \end{tabular}
\end{center}
For reference, the elementary Prange iteration counts
$\log_2\bigl(\binom{N}{k}/\binom{N-t}{k}\bigr)$ for the six sets are
$121.3$ / $136.8$ / $187.0$ / $202.3$ / $252.2$ / $267.5$: the Stern
figures sit $3.5$--$6.8$ bits above the bare iteration count, which
charges nothing per iteration.
Restoring ISD at $w=n$ costs $2.34$--$3.08\times$ in $N$ plus field bumps, for
$3.36$--$4.56\times$ public keys versus spec.

\subsection{GIJS at $w=n$ and clustering}\label{sec:wn-recovery}

The single-run GIJS distinguisher (one hold-out) at these parameters, with
$n'\leq\min(N,q)$ ($q^2$ via quadratic extension where $n$ sits close to $q$):
\begin{center}
    \begin{tabular}{llccc}
        \toprule
                 & field & best $(k',n',s,d)$ & $C_{\mathrm{run}}$ \\
        \midrule
        128-raw  & $q$   & (144,1742,3,4)     & $2^{72.4}$         \\
        128-cons & $q$   & (152,1956,3,4)     & $2^{73.4}$         \\
        192-raw  & $q^2$ & (177,2642,3,4)     & $2^{76.0}$         \\
        192-cons & $q^2$ & (184,2849,3,4)     & $2^{76.65}$        \\
        256-raw  & $q$   & (205,3529,3,4)     & $2^{78.5}$         \\
        256-cons & $q$   & (211,3737,3,4)     & $2^{79.0}$         \\
        \bottomrule
    \end{tabular}
\end{center}

The frequency-table attack of \S\ref{sec:phase1} needs pure columns; at
$w=n$ there are none, so that route is closed structurally. The
earlier version of this section recovered the pairing by counting
pair co-occurrences instead: a column is noisy in the shortened
code iff both it and its twin were unhit, so a true pair co-occurs
with probability $p_2 = \binom{N-2}{n'-2}/\binom{N}{n'}$, while a
non-pair needs four columns unhit, $p_4 =
\binom{N-4}{n'-4}/\binom{N}{n'}$, and the ratio $p_2/p_4 \approx
(N/n')^2$ is $1.54$--$1.63$ at the parameters above. That gap is
structural, not heuristic, but counting it out needed about $500$
trials to separate the smallest true count from the largest false
count with probability $0.99$ in the model.

Clustering (\S\ref{sec:cluster}) does the same job with far
fewer trials, and it is the natural tool at $w=n$, because there
is no pure set to confuse it: every column belongs to a pair, so
a clean clustering should return $w$ groups of size $2$ and an
empty all-zero group, with orphans and oversized groups only as
anomalies. The trial count is the clustering plan, not the
frequency plan. With $p = (n')_2/(N)_2$ the probability a given
pair survives a shortening, two distinct pairs fail to separate
in one trial with probability at most $p^2 + (1-p)^2$, and a pair
is invisible with probability $(1-p)^T$, so
$T = \max(T_{\mathrm{sep}}, T_{\mathrm{vis}})$ with
$\binom{w}{2}(p^2 + (1-p)^2)^T \le 0.01$ and
$w(1-p)^T \le 0.01$ gives, at the six sets above:

\begin{center}
    \begin{tabular}{lrrrr}
        \toprule
                 & $w$  & $N$  & $n'$ & $T_{\mathrm{cluster}}$ \\
        \midrule
        128-raw  & 1080 & 2160 & 1742 & 30 \\
        128-cons & 1218 & 2436 & 1956 & 30 \\
        192-raw  & 1665 & 3330 & 2642 & 30 \\
        192-cons & 1801 & 3602 & 2849 & 30 \\
        256-raw  & 2246 & 4492 & 3529 & 31 \\
        256-cons & 2382 & 4764 & 3737 & 31 \\
        \bottomrule
    \end{tabular}
\end{center}

That is $30$--$31$ trials, not about $500$. The same caveat
applies as at the shipped sets: clustering rests on exact
flagging (H5), and exact flagging at these much larger
$(k', n')$ has not been demonstrated; the trial counts above
are the plan \emph{if} flagging is exact, not a measured
result.

Once the pairing is known, the old per-pair disentangling ---
trying every $\alpha \in \Fq^*$ for every pair and testing each
candidate with GIJS over several shortenings --- is no longer
needed, and its cost should not be added. Instead, unmix three
pairs by the Phase~2 root-finding of \S\ref{sec:fastp2}, run
Sidelnikov--Shestakov to recover the curve, and then unmix the
remaining pairs directly against that curve, exactly as at the
shipped sets. The disentangling step is therefore polynomial
in $k$ and $N$, not a GIJS search per pair, and the cost of
the $w=n$ attack is, to first order, the cost of the clustering
trials themselves. What that cost is in field operations and
in wall time at these parameters, where $(n'-1)J$ sits close
to $N_d$ and reuse memory has not been measured, is left to
\S\ref{sec:total-costs}; the old total below is not the new
total.

\textbf{Old full costs, kept for comparison only}
(pair-frequency $+$ disentangling, superseded as above):
\begin{center}
    \begin{tabular}{lcccc}
        \toprule
                 & $T$ & pair-freq   & disentangle & total                \\
        \midrule
        128-raw  & 501 & $2^{92.2}$  & $2^{93.5}$  & $\mathbf{2^{94.0}}$  \\
        128-cons & 501 & $2^{93.3}$  & $2^{94.6}$  & $\mathbf{2^{95.1}}$  \\
        192-raw  & 496 & $2^{96.3}$  & $2^{97.7}$  & $\mathbf{2^{98.2}}$  \\
        192-cons & 496 & $2^{97.1}$  & $2^{98.5}$  & $\mathbf{2^{98.9}}$  \\
        256-raw  & 500 & $2^{99.3}$  & $2^{101.6}$ & $\mathbf{2^{101.9}}$ \\
        256-cons & 501 & $2^{99.8}$  & $2^{102.2}$ & $\mathbf{2^{102.5}}$ \\
        \bottomrule
    \end{tabular}
\end{center}
All six fell $33$--$154$ bits below target in that old model
(model estimates, as in Section~\ref{sec:key-cost}). The old
total is an overestimate on two counts: it pays about $500$
trials where clustering plans $30$--$31$, and it adds a
per-pair GIJS disentangling cost that Phase~2 root-finding
replaces with a polynomial step. The new total is therefore
lower than the old total, but it is not written here, because
the per-trial cost and memory at these $(k', n')$ have not
been measured on the new machine; see the template in
\S\ref{sec:total-costs}. The $w=n$ defense ---
$3.36$--$4.56\times$
keys to eliminate pure columns --- is futile on the mechanism
alone: clustering recovers the pairing without any
pure-column bootstrap, and Phase~2 finishes from three pairs.

\bibliographystyle{alpha}
\bibliography{references}

\section*{AI Usage Statement}
Muse AI (Meta) was used in the preparation of this paper for proofs,
calculations, and writing, and the creation of the faster non-trivial Phase 2 as well as, with careful 
guidance from the author, the reuse idea to speed up Phase 1. 

The other key ideas and insights, including that
the Ghoshal--Ishai--Jain--Sun distinguisher can tolerate  noisy rows, the shortening and 
counting appearances of twin pairs idea for twin pair recovery are the author's.

The author takes full responsibility for the
contents.

\end{document}
